\section*{Bab \ref{ch:g3:numbers} --- Bilangan sampai 10\,000}

\begin{solution}{exo:g3:numbers:1}
$2\,345$; \quad $6\,050$; \quad $9\,001$.
\end{solution}

\begin{solution}{exo:g3:numbers:2}
$1\,208$: seribu dua ratus delapan. \quad
$4\,070$: empat ribu tujuh puluh. \quad
$9\,999$: sembilan ribu sembilan ratus sembilan puluh sembilan.
\end{solution}

\begin{solution}{exo:g3:numbers:3}
Angka puluhan: $8$. Angka ribuan: $7$. Ratusan yang dimuat
seluruhnya: $73$ (karena $7\,382 = 73$ ratusan $+ 82$, yaitu
$73 \times 100 + 82$).
\end{solution}

\begin{solution}{exo:g3:numbers:4}
$5\,631 = 5$ ribuan $+ 6$ ratusan $+ 3$ puluhan $+ 1$ satuan.

$2\,046 = 2$ ribuan $+ 0$ ratusan $+ 4$ puluhan $+ 6$ satuan.

$8\,500 = 8$ ribuan $+ 5$ ratusan.
\end{solution}

\begin{solution}{exo:g3:numbers:5}
$456 < 465$ (puluhan: $5 < 6$); \quad
$1\,203 > 989$ (empat angka mengalahkan tiga); \quad
$6\,090 > 6\,009$ (puluhan: $9 > 0$); \quad
$9\,999 < 10\,000$.
\end{solution}

\begin{solution}{exo:g3:numbers:6}
$353 < 503 < 530 < 3\,500 < 5\,033$.
\end{solution}

\begin{solution}{exo:g3:numbers:7}
$150$ terletak di tengah antara $100$ dan $200$; $480$ tepat sebelum
$500$; $820$ tepat setelah $800$; $990$ hampir mencapai $1\,000$.
\end{solution}

\begin{solution}{exo:g3:numbers:8}
\omterm{def:g3:numbers:evenodd}{Genap}: $34$, $70$, $6\,548$ (angka terakhir $4$, $0$, $8$). \omterm{def:g3:numbers:evenodd}{Ganjil}:
$57$, $2\,481$ (angka terakhir $7$, $1$).
\end{solution}

\begin{solution}{exo:g3:numbers:9}
Dari $2\,970$: $2\,980$, $2\,990$, $3\,000$, $3\,010$, $3\,020$.

Dari $860$: $960$, $1\,060$, $1\,160$, $1\,260$.

Dari $9\,996$: $9\,997$, $9\,998$, $9\,999$, $10\,000$ --- langkah
terakhir menyeberang ke bilangan lima angka: batas seluruh bab kita!
\end{solution}

\begin{solution}{exo:g3:numbers:10}
Terbesar: $7\,530$ (angkanya berurutan mengecil). Terkecil: $3\,057$
(angka bukan \omterm{def:g1:counting:zero}{nol} yang terkecil di depan, lalu sisanya membesar ---
angka $0$ tepat setelah $3$).
\end{solution}

\begin{solution}{exo:g3:numbers:11}
Lebih besar daripada $800$: angka ratusannya $8$. \omterm{def:g3:numbers:evenodd}{Genap}: angka
satuannya harus $2$ (memakai $2$, $5$, $8$ masing-masing sekali).
Jadi bilangannya $852$.

\end{solution}
